%%%%%%%%%%%%%%%%%%%%%%%%%%%%%%%%%%%%%%%%%%%%%%%%%%%%%%%%%%%%%%%%%%%%%%%%%%%%%%%%
% program name style
\newcommand{\prognameplain}{Diabat}

\newcommand{\proglinewidth}{0.5pt}

\newcommand{\progchar}[2]{%
  \begingroup
  \colorlet{progfill}{white!#2!black}%
  \pdfrender{
    TextRenderingMode=FillStroke,
    LineWidth=\proglinewidth,
    FillColor=progfill,
    StrokeColor=black
  }{#1}%
  \endgroup
}

\DeclareRobustCommand{\progname}{%
  \scalebox{0.85}{%
    {\fontfamily{qcs}\fontseries{m}\selectfont
      \progchar{D}{50}%
      i%
      a%
      \progchar{b}{50}%
      a%
      \progchar{t}{50}%
    }%
  }%
}

\newcommand{\exec}{\texttt{diabat}}
\newcommand{\Sec}[1]{Sec.\,\ref{#1}}

\newcommand{\manualnote}[1]{%
  {#1}%
}

% key text to emphasize
\definecolor{keyemphgray}{gray}{0.25}
\newcommand{\keyemph}[1]{%
  \textcolor{keyemphgray}{\emph{#1}}%
}

%%%%%%%%%%%%%%%%%%%%%%%%%%%%%%%%%%%%%%%%%%%%%%%%%%%%%%%%%%%%%%%%%%%%%%%%%%%%%%%%
% Keyword description block
%%%%%%%%%%%%%%%%%%%%%%%%%%%%%%%%%%%%%%%%%%%%%%%%%%%%%%%%%%%%%%%%%%%%%%%%%%%%%%%%
\newcommand{\OpTrue}{{.true.}}
\newcommand{\OpFalse}{{.false.}}
\newcommand{\OpInt}{INTEGER}
\newcommand{\OpReal}{REAL}
\newcommand{\OpLogic}{LOGICAL}
\newcommand{\OpString}{STRING}
\newcommand{\OpIntseq}{INTEGER SEQUENCE}
\newcommand{\OpRealseq}{REAL SEQUENCE}
\newcommand{\OpNone}{NONE}
\newcommand{\OpOptional}{OPTIONAL}
\newcommand{\OpRequired}{REQUIRED}

\definecolor{keywordcolor}{rgb}{0.2,0.2,0.2}
\definecolor{backgray}{gray}{0.95}

\newcommand{\keyref}[2]{%
  {\hypersetup{linkcolor=black}\hyperlink{#1}{\texttt{#2}}}
}


% ¶¨ÒåÒ»¸öÐÂÃüÁîÀ´´´½¨Ò»¸ö¹Ø¼ü´ÊÃèÊö
\newcommand{\keyword}[9]{
    {
    \renewcommand{\arraystretch}{1.35}% ÉèÖÃÐмä¾à
    \fontsize{10pt}{10pt}\selectfont % ÉèÖÃ×ÖÌå´óСºÍÐоà
    \fontfamily{ptm}\selectfont % ÉèÖÃ×ÖÌåΪ Times
    \begin{longtable}{p{0.12\textwidth} p{0.78\textwidth}}
        \rowcolor{backgray}\textcolor{keywordcolor}{\textbf{Identifier}} & \raisebox{0pt}[0pt][0pt]{\hypertarget{#1}{\texttt{#2}}} \\
        \textcolor{keywordcolor}{\textbf{Description}} & #3 \\
        \textcolor{keywordcolor}{\textbf{Necessity}} & #4 \\
        \textcolor{keywordcolor}{\textbf{Data Type}} & #5 \\
        \textcolor{keywordcolor}{\textbf{Default}} & #6 \\
        \textcolor{keywordcolor}{\textbf{Options}} &  #7 \\
        \textcolor{keywordcolor}{\textbf{Suggestion}} & #8  \\
        \textcolor{keywordcolor}{\textbf{Example}} & \texttt{#9}
    \end{longtable}
    \vspace{-0.2cm}
    }
}



%%%%%%%%%%%%%%%%%%%%%%%%%%%%%%%%%%%%%%%%%%%%%%%%%%%%%%%%%%%%%%%%%%%%%%%%%%%%%%%%
% Semantic aliases for future chapters. Existing macros are preserved exactly.
%%%%%%%%%%%%%%%%%%%%%%%%%%%%%%%%%%%%%%%%%%%%%%%%%%%%%%%%%%%%%%%%%%%%%%%%%%%%%%%%
\let\ProgramName\progname
\let\ExecutableName\exec
\let\KeywordEntry\keyword
\let\ScriptCaption\scriptnum
\let\ExampleCaption\examplenum

\newcommand{\InternalNote}[1]{%
  \ifinternalmanual
  \begin{quote}
  \textbf{Internal note.} #1
  \end{quote}
  \fi
}

% Clickable keyword overview using the original keyword table hyperlink targets.
\newenvironment{keywordoverview}{%
  \begin{longtable}{p{0.34\textwidth}p{0.54\textwidth}}
  \rowcolor{backgray}\textbf{Keyword} & \textbf{Purpose}\\\hline
}{\end{longtable}}
\newcommand{\keyitem}[3]{\keyref{#1}{#2} & #3\\}

% Output filename entries are numbered within each section.  The heading and
% its description share the same indentation and cannot be separated at a page
% boundary.  Typography uses the existing body size and a bold mono filename.
\newcounter{outputfile}[section]
\newcommand{\outputfile}[2]{%
  \refstepcounter{outputfile}%
  \par\addvspace{0.48\baselineskip}%
  \begin{list}{\textcolor{keywordcolor}{\arabic{outputfile}.}}{%
    \setlength{\leftmargin}{2.05em}%
    \setlength{\labelwidth}{1.35em}%
    \setlength{\labelsep}{0.7em}%
    \setlength{\topsep}{0pt}%
    \setlength{\itemsep}{0pt}%
    \setlength{\parsep}{0pt}%
    \setlength{\partopsep}{0pt}%
  }
    \item \begin{minipage}[t]{\linewidth}%
      \noindent{\ttfamily\bfseries\detokenize{#1}}\par\nobreak
      \vspace{0.06\baselineskip}%
      \noindent #2\par
    \end{minipage}%
  \end{list}\par\addvspace{0.38\baselineskip}%
}
% A consistent unindented heading precedes each reproduced input listing.
\newcommand{\exampleheading}[1]{%
  \par\addvspace{0.35\baselineskip}%
  \noindent\textbf{#1}\par\nobreak
}
\newcommand{\examplepath}[1]{%
  \noindent{\footnotesize\path{#1}}\par\nobreak\smallskip
}
